\documentclass[11pt]{article}
\usepackage[margin=1in]{geometry}
\usepackage{amsmath,amssymb,amsthm,mathrsfs,hyperref}
\newtheorem{lemma}{Lemma}\newtheorem{theorem}[lemma]{Theorem}
\newtheorem{corollary}[lemma]{Corollary}
\title{The physical raising calculus at every prime}
\author{Internal research note}\date{18 November 2026}
\begin{document}\setlength{\emergencystretch}{1em}\maketitle
Let $p$ be any prime, $q=p-1$, and $N\ge1$.
Put$C=N+q-1$, $A=C-zD_z+z$,
$\Pi_N(X)=\prod_{\ell=0}^{N-1}(X+q\ell)$.
For$Q\in\mathbb Q_p[X]$ use the finite polynomial identity
\[
 \frac{\Pi_N(A)Q(A)1}{q^NN!}=\sum_jb_{N,Q,j}z^j,\qquad
 V_{N,a}(Q;v)=a!\sum_{\ell\ge0}p^\ell b_{N,Q,a+q\ell}v^\ell,
                 \quad0\le a<q .
\]
Write$s=s_p$ for the base-$p$ digit sum, and
\[
 e_a=\left\lceil\frac{s(N)-a-1}{q}\right\rceil\ge0,\qquad
                  \widetilde V_a=p^{-e_a}V_{N,a}.
\]

\begin{lemma}[Actual all-index digit content]\label{lem:content}
For every$Q\in\mathbb Z_p[X]$,
$V_{N,a}(Q)\in p^{e_a}\mathbb Z_p[v]$.
The same conclusion holds for integral coefficient polynomials
over any finite extension, with the same integer exponent$e_a$.
\end{lemma}
\begin{proof}
For integral series$f(t,v)$ define$f\in\mathcal F_j$ if every
coefficient of$t^n$ has Gauss valuation at least$(s(n)-j)/q$.
Then
\[
 \mathcal F_i\mathcal F_j\subseteq\mathcal F_{i+j},\qquad
 p\mathcal F_j\subseteq\mathcal F_{j-q},\qquad
 \mathcal F_i\subseteq\mathcal F_j\ (i\le j).
\]
Among integral series$\mathcal F_1$ is closed under$p$th powers.
A mixed$p$th-power multinomial has a factor$p$ and valuation
at least$1+\sum(s(n_i)-1)/q\ge(s(\sum n_i)-1)/q$.
Pure terms use$s(pn)=s(n)$ and integrality.

The exact diagonal series, obtained by formal Lagrange
substitution, is
\[
 \sum_{n\ge0}\frac{B_{n,n,q}(z)}{q^nn!}t^n
     =\kappa(t)e^{zH(t)},\quad
 R-R^p=t,\ R(0)=1,\ H=1-R,\quad
 \kappa=\frac{qR}{pR^q-1}.
\]
Here $B_{n,\gamma,q}=\Pi_n(\gamma+q-1-zD_z+z)1$. The fixed-shape series is
$(1-u)^{-(\gamma+q-1)/q}
 e^{z(1-(1-u)^{1/q})}$.
For $\gamma=n$, diagonal Lagrange with $\alpha=(1-u)^{-1/q}$
and $R=(1-u)^{1/q}$ yields
the stated numerator and denominator.

The integral series$H$ solves
\[
 H=\frac tq+\sum_{j=2}^p
                 \frac{(-1)^j}{q}\binom pj H^j.
\]
The middle coefficients are$p$-divisible, and the top term is
a$p$th power. Iteration is$t$-adically contracting and preserves
$\mathcal F_1$, so$H\in\mathcal F_1$.
This includes$p=2$, where$H=t+H^2$.

Use the formal Artin--Hasse integrality statement in
\texttt{lem:artin-hasse-closure} of the
\href{ARITHMETIC_FRONTIER_TERMINAL_ORDERS.tex}
{arithmetic note}.
For$E_i=(p^i-1)/q$, the equation
\[
 \mathscr A+\sum_{i\ge1}p^{E_i-i}v^{E_i}\mathscr A^{p^i}=H
\]
has an integral formal solution in$t\mathbb Z_p[v][[t]]$,
because$E_i-i\ge0$ and every nonlinear term has degree at least$p$.
Iteration and$p$th-power closure give$\mathscr A\in\mathcal F_1$.
At$z^q=pv$, formally$e^{zH}=E_p(z\mathscr A)$.
Every fixed$t$-coefficient is a finite polynomial in$z$:
no exponential is evaluated at a boundary point.
If$E_p(X)=\sum e_jX^j$, its normalized$a$th cyclic section is
\[
 a!\sum_{\ell\ge0}e_{a+q\ell}p^\ell v^\ell
                         \mathscr A^{a+q\ell}\in\mathcal F_a.
\]
Also
$\kappa=-q\sum_{j\ge0}p^j(1-H)^{qj+1}\in\mathcal F_1$,
coefficientwise$p$-adically convergently:
each term$p^jH^k$ has$k-qj\le1$.
The product proves the asserted content for$Q=1$.

The exact raising identities are
\[
 \begin{aligned}
 V_a(XQ)&=(C-a-qvD_v)V_a(Q)+aV_{a-1}(Q),\quad 1\le a<q,\\
 V_0(XQ)&=(C-qvD_v)V_0(Q)+\frac{pv}{(q-1)!}V_{q-1}(Q).
 \end{aligned}
\]
Here$e_{a-1}\ge e_a$ and$1+e_{q-1}\ge e_0$.
They preserve the content and prove it for all integral$Q$.
The extension-field assertion follows by finite linearity.
\end{proof}

Define exact operators on$q$-component columns by
\[
 D=\operatorname{diag}(C-a-qvD_v),\qquad L=D+J,
\]
\[
 J_{a,a-1}=a\,p^{e_{a-1}-e_a}\ (a\ge1),\qquad
 J_{0,q-1}=\frac{p^{1+e_{q-1}-e_0}v}{(q-1)!}.
\]
Then$\widetilde V(XQ)=L\widetilde V(Q)$.
For$p=2$ these are one-by-one operators$D=N-vD_v$, $J=2v$.

\begin{lemma}[A formal integral gauge]\label{lem:gauge}
One has$J^q=pvI$ and$[D,J]=-J$.
The formal matrices
\[
 E_\pm=\sum_{j\ge0}\frac{(\pm J)^j}{j!}
                    \in\operatorname{GL}_q(\mathbb Z_p[[v]])
\]
are inverse and$L=E_+DE_-$.
Their coefficients fail to tend to zero for every prime $p$.
\end{lemma}
\begin{proof}
The link exponents of$p$ are zero or one and sum to one.
The internal constants multiply to$(q-1)!$, cancelling the wrap
denominator. Hence$J^q=pvI$.
An internal diagonal difference gives$-1$; the wrap difference
$q-1$ and Euler contribution$-q$ also sum to$-1$.

For$j=q\ell+r$, $0\le r<q$, each nonzero matrix entry of$J^j$
has valuation$\ell+\epsilon$, $\epsilon\in\{0,1\}$.
Legendre's formula gives
\[
 v_p(J^j/j!)=\epsilon+\frac{s_p(j)-r}{q}\ge0.
\]
The digit sum is congruent to$r$ modulo$q$ and at least$r$;
for$r=0,j>0$ it is positive.
Each fixed$v$-coefficient receives only finitely many terms.
The constant matrices are triangular unipotent, so the two
formal series are invertible.
The exponential product and conjugation identities hold
coefficientwise, using$[J,D]=J$.

For$p>2$ choose any link of valuation zero and$j=p^m$.
It contributes a unit coefficient at arbitrarily high degree,
with no other$j$ contributing to that same entry and degree.
For$p=2$, the scalar coefficient$(2v)^j/j!$ has valuation
$s_2(j)=1$ along$j=2^m$.
In both cases the coefficients fail to tend to zero.
The identities are coefficientwise formal identities.
\end{proof}

\begin{theorem}[Contractive integer-valued calculus]\label{thm:int}
If$Q\in\mathbb Q_p[X]$ is integer-valued on$\mathbb Z_p$,
then$Q(L)$ preserves$\mathbb Z_p[v]^q$ and
$\mathbb Z_p\langle v\rangle^q$. On the Tate algebra
\[
 \|Q(L)\|\le\sup_{x\in\mathbb Z_p}|Q(x)|_p.
\]
Consequently the actual content assertion of
Lemma~\ref{lem:content} holds for every such$Q$, even with
rational polynomial coefficients.
\end{theorem}
\begin{proof}
On row$a$, degree$k$, the operator$\binom Dm$ multiplies by
$\binom{C-a-qk}{m}\in\mathbb Z_p$.
Thus it preserves integral formal series.
The formal gauge gives
$\binom Lm=E_+\binom Dm E_-$ with the same property.
On polynomial vectors this is a finite differential-polynomial
operator, so its output is both polynomial and integral.

The finite Newton expansion
$Q(X)=\sum_m b_m\binom Xm$ has
$b_m=\Delta^mQ(0)\in\mathbb Z_p$, and
$\max_m|b_m|_p=\sup_{\mathbb Z_p}|Q|_p$.
One inequality follows from differences of values at integers;
the other follows from the binomial expansion on$\mathbb Z_p$.
It gives the norm bound on polynomial vectors.
Their Gauss-norm density extends it to the Tate algebra.
Apply$Q(L)$ to the actual integral seed$\widetilde V(1)$.
This argument evaluates no formal gauge at$|v|_p=1$.
\end{proof}

\begin{corollary}[Additional divisibility]\label{cor:extra-content}
For $m\ge0$ and every integer-valued $Q\in\mathbb Q_p[X]$ on $\mathbb Z_p$,
\[
 V_{N,a}((X^p-X)^m Q;v)\in p^{e_a+m}\mathbb Z_p[v].
\]
\end{corollary}
\begin{proof}
By Fermat's congruence, $(X^p-X)/p$ is integer-valued on $\mathbb Z_p$.
Apply Theorem~\ref{thm:int} to its $m$th power times $Q$, then
multiply by $p^m$.
\end{proof}

\begin{corollary}[Continuous calculus and the residue tree]
\label{cor:continuous}
The polynomial calculus extends uniquely to a contractive
unital algebra homomorphism
\[
 C(\mathbb Z_p,\mathbb Q_p)\longrightarrow
       \mathcal B(\mathbb Q_p\langle v\rangle^q),\qquad f\mapsto f(L).
\]
The characteristic functions of residue classes give bounded
orthogonal projectors$P_{a,r}$, summing to one and satisfying
$P_{a,r}=\sum_{b\bmod p^{r+1},\,b\equiv a\bmod p^r}P_{b,r+1}$.
\end{corollary}
\begin{proof}
For$a=\sum_{i<r}a_ip^i$ set
\[
 \rho_{a,r}(X)=\prod_{i=0}^{r-1}
       \left[1-\left(\binom X{p^i}-a_i\right)^{p-1}\right].
\]
Lucas' identity and continuity show that$\rho$ is1 modulo$p$
on that residue class and0 elsewhere.
The polynomials$\rho^{p^m}$ converge uniformly to its indicator:
on the class$v_p(\rho^{p^m}-1)\ge m+1$; off it
$v_p(\rho^{p^m})\ge p^m$.
These bounds also hold at$p=2$.
Finite linear combinations approximate every locally constant
function; such functions are uniformly dense in
$C(\mathbb Z_p,\mathbb Q_p)$ by compactness.
Thus polynomials are dense, with an explicit approximation.

The preceding theorem applies after scalar scaling to every
polynomial, and bounds the operator norm by its uniform scalar norm.
Completeness gives the continuous extension; multiplication
and the projector identities pass to operator-norm limits.
Evaluated projector components of the actual seed can cancel.
An exclusion of their simultaneous cancellation remains open.
\end{proof}
\end{document}
